\documentclass{article}


% if you need to pass options to natbib, use, e.g.:
%     \PassOptionsToPackage{numbers, compress}{natbib}
% before loading neurips_2025


% ready for submission
\usepackage[preprint]{neurips_2025}


% to compile a preprint version, e.g., for submission to arXiv, add add the
% [preprint] option:
%     \usepackage[preprint]{neurips_2025}


% to compile a camera-ready version, add the [final] option, e.g.:
%     \usepackage[final]{neurips_2025}


% to avoid loading the natbib package, add option nonatbib:
%    \usepackage[nonatbib]{neurips_2025}


\usepackage[utf8]{inputenc} % allow utf-8 input
\usepackage[T1]{fontenc}    % use 8-bit T1 fonts
\usepackage{hyperref}       % hyperlinks
\usepackage{url}            % simple URL typesetting
\usepackage{booktabs}       % professional-quality tables
\usepackage{amsmath}        % Advanced math environments and commands
\usepackage{amssymb}        % Additional mathematical symbols
\usepackage{amsfonts}       % blackboard math symbols
\usepackage{mathtools}      % Extends amsmath with additional features
\usepackage{nicefrac}       % compact symbols for 1/2, etc.
\usepackage{microtype}      % microtypography
\usepackage{xcolor}         % colors

% Algorithm packages
\usepackage{algorithm}       % For algorithm environment
\usepackage{algpseudocode}   % For algorithmic environment


\title{DCIDM: Dual Class and Instance-Guided Diffusion for Targeted Synthetic Data Generation}


% The \author macro works with any number of authors. There are two commands
% used to separate the names and addresses of multiple authors: \And and \AND.
%
% Using \And between authors leaves it to LaTeX to determine where to break the
% lines. Using \AND forces a line break at that point. So, if LaTeX puts 3 of 4
% authors names on the first line, and the last on the second line, try using
% \AND instead of \And before the third author name.


\author{%
  Mykhailo~Sakevych \\
  Department of Computer Science\\
  Texas State University\\
  San Marcos, TX 78666 \\
  \texttt{ukb12@txstate.edu} \\
  % examples of more authors
  \And
  Vangelis~Metsis \\
  Department of Computer Science\\
  Texas State University\\
  San Marcos, TX 78666 \\
  \texttt{vmetsis@txstate.edu} \\
}


\begin{document}


\maketitle


\begin{abstract}
Class imbalance severely hampers machine learning model performance, particularly on minority classes. This is acute for time series data, where traditional oversampling methods often fail to capture complex temporal dynamics or generate diverse, high-fidelity samples. We introduce DCIDM, a novel Dual Conditioned Instance-guided Diffusion Model, to generate targeted synthetic minority class samples. DCIDM uniquely combines classifier-free guidance for robust class conditioning with a novel instance-level guidance strategy. We first systematically identify sparse, boundary-critical regions in a learned feature embedding space where additional minority samples would be most beneficial. Then, at inference, the reverse diffusion process is initiated from intermediate timesteps using an aggregation of partially noised instances from these identified target regions, effectively steering generation towards areas crucial for improving classifier decision boundaries without requiring model retraining. Comprehensive evaluations on the highly imbalanced datasets demonstrate that DCIDM significantly improves classification performance, particularly for underrepresented classes, outperforming existing augmentation techniques and mitigating model collapse. Our approach offers a principled framework for enhancing model robustness and generalization on imbalanced time series datasets.
\end{abstract}


\section{Introduction}

Class imbalance represents a persistent and critical challenge in machine learning, where the preponderance of majority class samples biases models, leading to poor performance on underrepresented minority classes \cite{Prati}. This issue is especially pronounced in high-stakes domains such as medical diagnostics (e.g., rare disease detection), fraud detection, and industrial anomaly detection, where minority classes often correspond to the most critical events. For time series data, common in healthcare monitoring (e.g., detecting rare cardiac arrhythmias) and industrial systems (e.g., predicting equipment failures), the problem is compounded by complex temporal dependencies and high dimensionality, making effective minority class representation even more difficult. This inherent imbalance biases learning towards majority classes, resulting in suboptimal decision boundaries and poor generalization for the crucial, yet scarce, minority instances.

\subsection{Limitations of Existing Approaches}
Conventional methods to address class imbalance include resampling techniques (undersampling majority, oversampling minority), algorithmic modifications, and synthetic data generation. Traditional oversampling methods like SMOTE and its variants (Borderline-SMOTE, ADASYN) attempt to create new minority samples by interpolating between existing ones \cite{Chawla_2002,10.1007/11538059_91, 4633969}. However, for time series, these methods often struggle to capture complex temporal dynamics, generating unrealistic samples that may not truly represent the minority class distribution or can even miscalibrate models \cite{carriero2024harmsclassimbalancecorrections}. While Generative Adversarial Networks (GANs) have been explored for time series augmentation \cite{li2022tts}, they frequently suffer from training instability and mode collapse, especially when dealing with highly imbalanced datasets.

\subsection{Diffusion Models for Imbalanced Learning}
Denoising Diffusion Probabilistic Models (DDPMs) have emerged as a powerful alternative, offering exceptional generation quality and training stability \cite{Ho_2020_DDPM, Sohl_Dickstein_2015_Diffusion}. They operate via a forward process that incrementally adds noise and a reverse process that learns to denoise and recover the data distribution. This step-by-step denoising is well-suited for capturing complex temporal dependencies in time series \cite{rasul2021autoregressive, li2024biodiffusion}. Classifier-free guidance further enhances conditional generation by incorporating conditioning information directly into the diffusion process without needing an explicit classifier, improving sample quality and efficiency \cite{Ho_Salimans_2022_ClassifierFree}.

However, even with advancements in conditional diffusion models for time series, a critical gap remains: existing models typically generate synthetic samples from \textit{random} or broad areas of the learned feature space. While these samples might be class-consistent, such untargeted augmentation may not optimally improve decision boundaries for imbalanced classes and, in some cases, introducing large amounts of random synthetic data can be less effective or even detrimental to classifier performance compared to strategically placed samples.

\subsection{Our approach: Targeted Generation with DCIDM}
We introduce the Dual Conditioned Instance-guided Diffusion Model (DCIDM), a novel framework designed to generate synthetic minority class time series samples by strategically targeting regions most critical for classification. DCIDM uniquely combines two conditioning mechanisms:
\begin{enumerate}
    \item \textbf{Class Conditioning via Classifier-Free Guidance:} Ensures generated samples adhere to the target minority class characteristics.
    \item \textbf{Targeted Instance Conditioning via "Partially Diffused Conditioning":} A novel inference-time strategy that guides the generation process towards specific, beneficial regions of the feature space. This is achieved by initiating the reverse diffusion process from an aggregation of partially noised minority-class instances identified as boundary-critical.
\end{enumerate}
Our approach first involves a systematic methodology to identify sparse, boundary-critical regions in a learned feature embedding space where additional minority samples would be most beneficial for improving classifier decision boundaries. Then, at inference, DCIDM steers generation towards these vital areas without requiring model retraining for each target region. This targeted synthesis aims to directly address the weaknesses in classifier decision boundaries caused by data scarcity.

Our work makes the following contributions:
\begin{enumerate}
    \item A specialized class-conditional DDPM architecture tailored for fixed-length time series generation, capable of capturing diverse temporal patterns.
    \item A systematic methodology for identifying sparse, boundary-critical regions in the feature embedding space where synthetic minority data is most impactful.
    \item A novel inference-time instance-guidance technique, "partially diffused conditioning," that directs the diffusion generation process towards these identified target regions without retraining the diffusion model.
    \item Comprehensive evaluation demonstrating significant improvements in classification performance on imbalanced time series datasets, particularly for underrepresented classes, validated on a variety of datasets.
    \item Theoretical analysis connecting the diffusion process dynamics with decision boundary refinement, establishing a formal framework for understanding how targeted synthetic augmentation affects classifier performance.
\end{enumerate}

We demonstrate empirically that DCIDM significantly enhances the accuracy and precision of predictive models on imbalanced time series datasets, explicitly mitigating common issues like model collapse sometimes seen with other augmentation techniques, and leading to consistent performance gains.




\section{Related Work}

Addressing class imbalance is a critical challenge in machine learning, and synthetic data generation has emerged as a prominent strategy. Our work builds upon advancements in generative models, particularly diffusion models, and introduces novel techniques for targeted, conditional synthesis of time series data.

\subsection{Synthetic Data Generation for Imbalanced Learning}
Traditional methods for handling imbalanced datasets often involve resampling. SMOTE (Synthetic Minority Over-sampling Technique) and its variants like Borderline-SMOTE and ADASYN generate synthetic minority samples by interpolating between existing instances \cite{Chawla_2002, 10.1007/11538059_91, 4633969}. While widely used, these methods can struggle with high-dimensional data like time series, potentially creating samples that do not faithfully represent the true data distribution or blur class distinctions.
More recent approaches leverage advanced generative models. Generative Adversarial Networks (GANs) \cite{Goodfellow_2014_GAN} and Variational Autoencoders (VAEs) \cite{Kingma_2013_VAE} have been explored for synthetic data generation. While GANs can produce realistic samples, they often suffer from training instability and mode collapse, particularly with limited minority class data, a challenge also present in time series applications \cite{li2022tts}. VAEs, while more stable, may produce less diverse samples. Other techniques include adaptive synthetic sampling (ADASYN), focusing on harder-to-learn minority examples, and cluster-based oversampling. Some research has also investigated using Large Language Models (LLMs) for generating synthetic tabular data \cite{Borisov_2023_LLMTabular}.

\subsection{Diffusion Models for Data Generation}
Denoising Diffusion Probabilistic Models (DDPMs) \cite{Ho_2020_DDPM, Sohl_Dickstein_2015_Diffusion} have recently shown state-of-the-art performance in generating high-fidelity samples across various domains, including images, audio, and, more recently, time series \cite{li2024biodiffusion}. They operate by progressively adding noise to data and then learning to reverse this process. Their ability to capture complex data distributions and their training stability make them attractive for synthetic data generation, often outperforming GANs in sample quality and diversity. Diffusion models have also been explored for tabular data generation and imputation \cite{kotelnikov2023tabddpm}.

\subsection{Conditional Diffusion Models}
A key advantage of diffusion models is their amenability to conditioning, allowing the generation process to be guided by auxiliary information such as class labels or other data modalities. Classifier guidance involves training an additional classifier on noisy data to steer the diffusion process \cite{Dhariwal_Nichol_2021_ClassifierGuidance}. A more common and often more effective approach is classifier-free guidance, where the diffusion model is jointly trained on conditional and unconditional objectives, enabling guidance at inference time without an explicit classifier \cite{Ho_Salimans_2022_ClassifierFree}. This method significantly improves sample quality and adherence to conditioning signals. Our approach leverages classifier-free guidance for robust class conditioning.

\subsection{Diffusion Models for Time Series Data}
The application of diffusion models to time series data is a growing area. Models like TimeGrad \cite{rasul2021autoregressive} and Diffusion-TS \cite{Gong_2023_DiffusionTS} have adapted diffusion principles for time series forecasting and generation. More recently, BioDiffusion \cite{li2024biodiffusion} demonstrated a versatile diffusion model for biomedical signal synthesis. For medical time series like ECG data, diffusion models are explored for augmenting datasets to improve tasks such as arrhythmia classification \cite{adib2023synthetic}. While these models can generate high-quality class-conditional time series, they typically sample from the entirety of the learned conditional distribution. This untargeted generation may not be optimal for specifically addressing class imbalance, where augmenting particular sparse or boundary-critical regions is more beneficial.

\subsection{Targeted and Instance-Conditioned Generation}
For imbalanced learning, augmenting specific, critical regions of the data distribution can be more effective than uniform oversampling. The concept of instance-conditioned generation, where the generative process is guided by specific data instances, has been explored, primarily in the image domain (e.g., InstanceDiffusion \cite{Wang_2023_InstanceDiffusion}).
Our approach, DCIDM, introduces a novel "partially diffused conditioning" technique. This method uses an aggregation of partially noised minority class instances, identified from systematically determined boundary-critical regions, to initiate and guide the reverse diffusion process at inference time. This allows for targeted generation towards these beneficial regions without requiring model retraining for specific instance guidance, a key distinction from prior instance-conditioning methods.

\subsection{Double Conditioning in Diffusion Models}
The idea of using multiple conditions to guide diffusion models has appeared in various contexts. For example, "double conditioning" has been used for out-of-distribution detection in medical imaging (class information and latent image features) \cite{mishra2023dual}, and for sequential recommendation (user and item context) \cite{huang2024dual}.
Our DCIDM also employs a form of dual conditioning: it simultaneously leverages class labels (via classifier-free guidance) and targeted instance-level information (via our novel "partially diffused conditioning" derived from boundary-critical region identification). The novelty lies in the specific combination and the inference-time mechanism designed to strategically augment sparse, critical regions for imbalanced time series classification.

\subsection{Distinction of Our Approach}
DCIDM distinguishes itself by:
\begin{enumerate}
    \item \textbf{Strategic Targeting for Imbalance:} Unlike general time series diffusion models that generate samples broadly, DCIDM focuses on systematically identified boundary-critical and sparse minority regions. This targeted approach aims to directly improve classifier decision boundaries where they are weakest due to data scarcity.
    \item \textbf{Novel Instance-Level Guidance Mechanism:} The "partially diffused conditioning" provides a unique inference-time method to steer generation towards desired regions using an aggregate of exemplars, without retraining the core diffusion model. This is distinct from prior instance or multi-conditional approaches.
    \item \textbf{Synergistic Dual Conditioning:} The combination of robust classifier-free guidance for class adherence and our targeted instance-level guidance ensures that generated samples are both class-consistent and strategically placed to remediate imbalance effectively in time series data.
\end{enumerate}
By integrating these elements, DCIDM aims to generate high-quality, strategically valuable synthetic time series data that more effectively addresses class imbalance challenges compared to existing methods.



\section{Methodology}

This section describes our framework for guided synthetic data generation using conditional diffusion models, focusing on time series classification with imbalanced classes. We first outline the problem formulation, then detail the components of our approach, from diffusion model architecture to the targeted generation strategy.

\subsection{Problem Formulation}

Let $\mathcal{D} = \{(\mathbf{x}_i, y_i)\}_{i=1}^N$ be a time series dataset with $N$ samples, where each $\mathbf{x}_i \in \mathbb{R}^T$ represents a time series of length $T$, and $y_i \in \{1, 2, \ldots, C\}$ denotes the corresponding class label. We consider imbalanced datasets where the number of samples per class follows a skewed distribution, with minority classes having significantly fewer instances than majority classes.

Our objective is to generate synthetic time series samples $\hat{\mathbf{x}}$ for minority classes that:
\begin{enumerate}
    \item Maintain high fidelity to the true data distribution $p(\mathbf{x}|y)$
    \item Specifically target decision boundary regions to improve classification performance
    \item Diversify the representation of minority classes in sparsely populated regions
\end{enumerate}

\subsection{Class-Conditional Diffusion Model for Time Series}

\subsubsection{Diffusion Process Formulation}

We employ a class-conditional denoising diffusion probabilistic model (DDPM) as our generative framework. The diffusion process consists of a forward process that gradually adds Gaussian noise to the data, and a reverse process that learns to denoise the data conditioned on class labels.

\paragraph{Forward Process:} We define a Markov chain that gradually adds noise to the original time series $\mathbf{x}_0 = \mathbf{x}$ over $T$ timesteps:

\begin{equation}
q(\mathbf{x}_t | \mathbf{x}_{t-1}) = \mathcal{N}(\mathbf{x}_t; \sqrt{1 - \beta_t} \mathbf{x}_{t-1}, \beta_t \mathbf{I})
\end{equation}

where $\{\beta_t\}_{t=1}^T$ is a fixed variance schedule with $\beta_t \in (0, 1)$. This process has the useful property that we can sample $\mathbf{x}_t$ directly from $\mathbf{x}_0$ via:

\begin{equation}
q(\mathbf{x}_t | \mathbf{x}_0) = \mathcal{N}(\mathbf{x}_t; \sqrt{\bar{\alpha}_t} \mathbf{x}_0, (1 - \bar{\alpha}_t) \mathbf{I})
\end{equation}

where $\alpha_t = 1 - \beta_t$ and $\bar{\alpha}_t = \prod_{s=1}^t \alpha_s$. This allows efficient sampling using the reparameterization trick:

\begin{equation}
\mathbf{x}_t = \sqrt{\bar{\alpha}_t} \mathbf{x}_0 + \sqrt{1 - \bar{\alpha}_t} \boldsymbol{\epsilon}, \quad \boldsymbol{\epsilon} \sim \mathcal{N}(\mathbf{0}, \mathbf{I})
\end{equation}

\paragraph{Reverse Process:} We train a neural network $\boldsymbol{\epsilon}_\theta(\mathbf{x}_t, t, y)$ to predict the noise component given the noisy time series $\mathbf{x}_t$, timestep $t$, and class label $y$. The model is trained to minimize:

\begin{equation}
\mathcal{L} = \mathbb{E}_{t, \mathbf{x}_0, \boldsymbol{\epsilon}, y} \left[ \| \boldsymbol{\epsilon} - \boldsymbol{\epsilon}_\theta(\sqrt{\bar{\alpha}_t} \mathbf{x}_0 + \sqrt{1 - \bar{\alpha}_t} \boldsymbol{\epsilon}, t, y) \|_2^2 \right]
\end{equation}

During sampling, we reverse the diffusion process by iteratively denoising from Gaussian noise:

\begin{equation}
\mathbf{x}_{t-1} = \frac{1}{\sqrt{\alpha_t}} \left( \mathbf{x}_t - \frac{1-\alpha_t}{\sqrt{1-\bar{\alpha}_t}} \boldsymbol{\epsilon}_\theta(\mathbf{x}_t, t, y) \right) + \sigma_t \mathbf{z}, \quad \mathbf{z} \sim \mathcal{N}(\mathbf{0}, \mathbf{I})
\end{equation}

where $\sigma_t^2 = \beta_t$ and $\mathbf{z}$ is introduced to match the variance of the posterior distribution.

\subsubsection{Model Architecture}

We adapt the U-Net architecture for one-dimensional time series data with the following components:

\begin{enumerate}
    \item \textbf{Feature Extraction:} A series of downsampling blocks with 1D convolutions that process the time series signal:
    \begin{equation}
    h_i = \text{DoubleConv1D}(\text{MaxPool1D}(h_{i-1})) + \text{Emb}(t)
    \end{equation}

    \item \textbf{Self-Attention Mechanism:} At multiple resolutions to capture long-range dependencies in the temporal domain:
    \begin{equation}
    \text{Attention}(Q, K, V) = \text{softmax}\left(\frac{QK^T}{\sqrt{d_k}}\right)V
    \end{equation}

    \item \textbf{Conditional Embeddings:} Class labels $y$ are embedded and added to the time embedding:
    \begin{equation}
    \tau_t + \text{Emb}(y) \in \mathbb{R}^d
    \end{equation}
    where $\tau_t$ is the positional encoding of timestep $t$

    \item \textbf{Skip Connections:} To preserve temporal structure during upsampling:
    \begin{equation}
    h_i^{up} = \text{UpConv1D}(\text{Concat}(h_i^{down}, h_{i-1}^{up})) + \text{Emb}(t)
    \end{equation}
\end{enumerate}

The complete architecture is implemented as a conditional U-Net with alternating layers of 1D convolutions, self-attention, and time/class embeddings, optimized for capturing the temporal characteristics of fixed-length time series.

\subsection{Feature Embedding Space Construction}

To identify meaningful regions for targeted generation, we construct a feature embedding space that captures the semantic structure of the time series dataset.

\begin{enumerate}
    \item \textbf{Base Classifier Training:} We train a convolutional neural network $f_\phi$ on the original dataset $\mathcal{D}$ to classify time series into their respective classes:
    \begin{equation}
    \hat{y} = \text{argmax}(f_\phi(\mathbf{x}))
    \end{equation}

    \item \textbf{Embedding Extraction:} We extract embeddings $\mathbf{z}_i$ from the penultimate layer of the classifier for all training instances:
    \begin{equation}
    \mathbf{z}_i = E_\phi(\mathbf{x}_i) \in \mathbb{R}^d
    \end{equation}
    where $E_\phi$ is the embedding function derived from the classifier network $f_\phi$.

    \item \textbf{Embedding Normalization:} We normalize the embeddings to ensure consistent scaling:
    \begin{equation}
    \mathbf{z}_i' = \frac{\mathbf{z}_i - \boldsymbol{\mu}}{\boldsymbol{\sigma}}
    \end{equation}
    where $\boldsymbol{\mu}, \boldsymbol{\sigma}$ are the mean and standard deviation vectors computed over all embeddings.
\end{enumerate}

This embedding space provides a compact representation where the proximity between samples reflects their semantic similarity, enabling more targeted analysis of decision boundaries.


\subsection{Identifying Target Regions for Synthetic Data}
\label{sec:identifying_target_regions} % Label for this main section

To guide the generation of synthetic minority samples towards areas where they can be most impactful for improving classifier performance, we employ a multi-stage strategy. This strategy first identifies critical regions in the learned feature embedding space based on characteristics of the existing minority class instances. Then, for each identified critical region, we select a set of $L$ nearby instances from the \textit{entire dataset} (regardless of their original class) to serve as spatial anchors for our instance-guided diffusion. The rationale is that while the DCIDM's class conditioning (as described in Section~\ref{ssec:classifier_free_guidance}) ensures the generated sample belongs to the target minority class, the instance conditioning (detailed in Section~\ref{ssec:partially_diffused_conditioning}) primarily steers generation towards a specific locale in the feature space. Using local exemplars of any class for this spatial guidance can provide richer information about the manifold structure, especially if minority instances are extremely sparse in that critical region.

Our approach synergistically combines several criteria to define the initial critical regions based on minority class data:
\begin{enumerate}
    \item \textbf{Boundary Proximity:} We identify instances of minority classes that are critical to defining the decision boundaries with other classes.
    \item \textbf{Classifier Uncertainty:} We pinpoint minority class instances where a pre-trained classifier exhibits high uncertainty.
    \item \textbf{Data Sparsity:} Within these boundary-critical or uncertain minority areas, we prioritize regions that are sparsely populated by existing minority class samples.
\end{enumerate}
Candidate minority instances identified through these criteria are then clustered to define distinct, coherent target regions ($\mathcal{C}_r$). These regions, defined by minority characteristics, then guide the selection of local representative instances from the broader dataset.

\subsubsection{Support Vector-Based Boundary Region Identification}
\label{sssec:sv_boundary_identification}

To identify minority instances intrinsically linked to decision boundaries, we leverage Support Vector Machines (SVMs). For each minority class $c_i \in \mathcal{C}_{minority}$ and each contrasting class $c_j$ (typically a majority class or another relevant class):

\begin{enumerate}
    \item \textbf{Train Binary SVM Classifier:} We train a binary SVM classifier, equipped with a Radial Basis Function (RBF) kernel, on the embeddings $\mathbf{z}$. This SVM$_{ij}$ is specifically trained to distinguish between instances of class $c_i$ and class $c_j$:
    \begin{equation}
    \text{SVM}_{ij} \text{ is trained on } \{ (\mathbf{z}_k, y_k') \mid y_k' \in \{c_i, c_j\} \}.
    \end{equation}
    The RBF kernel allows the SVM to find non-linear decision boundaries in the embedding space.

    \item \textbf{Identify Support Vectors:} From the trained SVM$_{ij}$, we extract its set of support vectors, denoted $\mathcal{SV}_{ij}$. These are the data points from the training set that lie closest to the decision boundary and are pivotal in defining its precise location and orientation.

    \item \textbf{Select Minority Class Support Vectors as Boundary Candidates:} We form a set of boundary candidate instances $\mathcal{B}_{ij}$ by selecting those support vectors from $\mathcal{SV}_{ij}$ that originally belong to the target minority class $c_i$:
    \begin{equation}
    \mathcal{B}_{ij} = \{\mathbf{x}_k \mid \mathbf{z}_k = E_\phi(\mathbf{x}_k), \mathbf{z}_k \in \mathcal{SV}_{ij}, \text{ and } y_k = c_i\}.
    \end{equation}
    This selection includes SVs that are correctly classified by SVM$_{ij}$ as well as any SVs from class $c_i$ that SVM$_{ij}$ may have misclassified. The rationale is that all such minority class SVs represent instances that are inherently difficult for the SVM to separate from class $c_j$, thus marking them as critical points near or straddling the decision boundary in the embedding space.
\end{enumerate}
The complete set of boundary-candidate minority instances, $\mathcal{B}$, is formed by the union of candidates from all such pairwise comparisons: $\mathcal{B} = \bigcup_{i,j} \mathcal{B}_{ij}$. These instances highlight regions of high inter-class contention relevant to the minority classes.

\subsubsection{Uncertainty-Based Filtering}
\label{sssec:uncertainty_filtering}

Complementary to the structural boundary identification by SVMs, we also identify minority class instances where the base deep learning classifier $f_\phi$ exhibits uncertainty.

\begin{enumerate}
    \item \textbf{Compute Prediction Entropy:} For each minority class instance $\mathbf{x}_i$, we compute the entropy $H(\mathbf{x}_i)$ of its predicted probability distribution $P(y=c|\mathbf{x}_i)$ obtained from the softmax output of $f_\phi$:
    \begin{equation}
    H(\mathbf{x}_i) = -\sum_{c=1}^C P(y=c|\mathbf{x}_i) \log_2 P(y=c|\mathbf{x}_i).
    \end{equation}
    Higher entropy indicates greater uncertainty in the classifier's prediction.

    \item \textbf{Select High-Uncertainty Instances:} We select minority class instances whose prediction entropy exceeds a predefined threshold $\gamma$:
    \begin{equation}
    \mathcal{U} = \{\mathbf{x}_i \mid H(\mathbf{x}_i) > \gamma, y_i \in \mathcal{C}_{minority}\}.
    \end{equation}
    These instances in $\mathcal{U}$ represent minority points where the classifier lacks confidence.
\end{enumerate}

\subsubsection{Combined Candidate Selection, Sparsity Filtering, and Clustering to Define Critical Minority Regions}
\label{sssec:combined_selection_clustering}

The subsequent steps integrate the boundary-critical minority candidates ($\mathcal{B}$) and high-uncertainty minority candidates ($\mathcal{U}$), filter them for sparsity relative to their own class, and then cluster them to define actionable target zones based on these minority characteristics.

\begin{enumerate}
    \item \textbf{Estimate Local Density for Minority Candidates:} For each minority instance $\mathbf{x}_i$ in the combined set of candidates ($\mathcal{B} \cup \mathcal{U}$), we estimate its local data density. We use the average distance to its $k$ nearest neighbors (k-NN) \textit{within its own minority class} $y_i$ in the embedding space as an inverse measure of density (i.e., a measure of sparsity):
    \begin{equation}
    \rho(\mathbf{z}_i) = \frac{1}{k} \sum_{\mathbf{z}_j \in \mathcal{N}_k(\mathbf{z}_i, y_i)} \|\mathbf{z}_i - \mathbf{z}_j\|_2,
    \end{equation}
    where $\mathcal{N}_k(\mathbf{z}_i, y_i)$ are the $k$ nearest neighbors of $\mathbf{z}_i$ that also belong to class $y_i$. A larger $\rho(\mathbf{z}_i)$ indicates that $\mathbf{z}_i$ is in a sparser region of its class.

    \item \textbf{Form Final Minority Candidate Set $\mathcal{S}$:} We combine the boundary-critical and high-uncertainty minority instances, and then apply the sparsity filter. A minority instance $\mathbf{x}_i$ is included in the final candidate set $\mathcal{S}$ if it belongs to either $\mathcal{B}$ or $\mathcal{U}$, and its corresponding embedding $\mathbf{z}_i$ is located in a relatively sparse area of its own class (its sparsity score $\rho(E_\phi(\mathbf{x}_i))$ is greater than a threshold $\lambda$):
    \begin{equation}
    \mathcal{S} = \{\mathbf{x}_i \mid \mathbf{x}_i \in (\mathcal{B} \cup \mathcal{U}) \land \rho(E_\phi(\mathbf{x}_i)) > \lambda \land y_i \in \mathcal{C}_{minority}\}.
    \end{equation}
    This ensures that $\mathcal{S}$ consists of minority class instances that are critical (boundary/uncertain) and reside in sparsely populated areas of their own class distribution.

    \item \textbf{Cluster Minority Candidates to Define Critical Regions $\mathcal{C}_r$:} We apply the DBSCAN algorithm to the embeddings $\{\mathbf{z}_i \mid \mathbf{x}_i \in \mathcal{S}\}$ of these final minority candidate instances. DBSCAN groups these points into clusters $\mathcal{C}_1, \mathcal{C}_2, \ldots, \mathcal{C}_R$.
    \begin{equation}
    \mathcal{C}_1, \mathcal{C}_2, \ldots, \mathcal{C}_R = \text{DBSCAN}(\{E_\phi(\mathbf{x}_i) \mid \mathbf{x}_i \in \mathcal{S}\}).
    \end{equation}
    Each cluster $\mathcal{C}_r$ now represents a coherent region in the feature space identified as critical and sparse based on the characteristics of the \textit{minority class instances} it contains. Let $\mathbf{z}_{\text{centroid},r}$ be the centroid (mean embedding) of the minority instances within cluster $\mathcal{C}_r$.

    \item \textbf{Select Local Representative Instances $\mathcal{R}_r$ from Entire Dataset:} For each identified critical minority region $\mathcal{C}_r$ (represented by its centroid $\mathbf{z}_{\text{centroid},r}$), we now select a representative set of $L$ instances, $\mathcal{R}_r$, from the \textit{entire original dataset} $\mathcal{D} = \{(\mathbf{x}_k, y_k)\}_{k=1}^N$. These representatives are chosen as the $L$ instances from $\mathcal{D}$ whose embeddings $E_\phi(\mathbf{x}_k)$ are closest to the region's centroid $\mathbf{z}_{\text{centroid},r}$, regardless of their class $y_k$.
    \begin{equation}
    \mathcal{R}_r = \{\mathbf{x}_k \in \mathcal{D} \mid \text{rank}(\|E_\phi(\mathbf{x}_k) - \mathbf{z}_{\text{centroid},r}\|_2) \leq L \}.
    \end{equation}
    These sets $\mathcal{R}_r$ provide local, spatially relevant exemplars—potentially from various classes—that characterize the immediate neighborhood of the critical minority region. They serve as the conditioning instances for the "partially diffused conditioning" step (detailed in Section~\ref{ssec:partially_diffused_conditioning}), guiding the generation spatially, while the target minority class label $c$ (used with classifier-free guidance as per Section~\ref{ssec:classifier_free_guidance}) ensures the appropriate class characteristics are synthesized.
\end{enumerate}
This comprehensive methodology first pinpoints critical regions based on minority class data, and then uses these regions to select local exemplars from the full dataset to guide generation spatially. This aims to maximize the utility of augmented data by leveraging broad local manifold information while ensuring class-specific synthesis.


\subsection{Instance-Guided Diffusion Sampling}
\label{sec:instance_guided_sampling}

We introduce a novel approach to guide the diffusion model toward generating samples in specific target regions:

\subsubsection{Partially Diffused Conditioning}
\label{ssec:partially_diffused_conditioning}

For a target region represented by instances $\mathcal{R}_r$ from a minority class $c$:

\begin{enumerate}
    \item Select a set of $K$ conditioning instances $\{\mathbf{x}_i^c\}_{i=1}^K \subset \mathcal{R}_r$.

    \item Apply partial forward diffusion to each conditioning instance up to a specified timestep $t'$:
    \begin{equation}
    \mathbf{x}_{i,t'} = \sqrt{\bar{\alpha}_{t'}} \mathbf{x}_i^c + \sqrt{1 - \bar{\alpha}_{t'}} \boldsymbol{\epsilon}_i
    \end{equation}
    where $\boldsymbol{\epsilon}_i \sim \mathcal{N}(\mathbf{0}, \mathbf{I})$ and $t'$ controls how much information from the original instance is preserved.

    \item Average the partially diffused instances to create a conditioning point:
    \begin{equation}
    \bar{\mathbf{x}}_{t'} = \frac{1}{K} \sum_{i=1}^K \mathbf{x}_{i,t'}
    \end{equation}

    \item Use this as the starting point for the reverse diffusion process:
    \begin{equation}
    \mathbf{x}_{t'} = \bar{\mathbf{x}}_{t'}
    \end{equation}

    \item Run the class-conditional reverse diffusion process from timestep $t'$ to $0$:
    \begin{equation}
    \mathbf{x}_{t-1} = \frac{1}{\sqrt{\alpha_t}} \left( \mathbf{x}_t - \frac{1-\alpha_t}{\sqrt{1-\bar{\alpha}_t}} \boldsymbol{\epsilon}_\theta(\mathbf{x}_t, t, c) \right) + \sigma_t \mathbf{z}
    \end{equation}
    for $t = t', t'-1, \ldots, 1$.
\end{enumerate}

This approach, which we call "partially diffused conditioning," guides the generation process toward creating realistic samples that maintain some characteristics of the boundary region while ensuring diversity through the stochastic nature of the diffusion process.

\subsubsection{Classifier-Free Guidance}
\label{ssec:classifier_free_guidance}

To further strengthen the conditioning on both class and instance information, we incorporate classifier-free guidance:

\begin{enumerate}
    \item At each denoising step, compute both conditional and unconditional noise predictions:
    \begin{equation}
    \boldsymbol{\epsilon}_\theta^c(\mathbf{x}_t, t, c) \quad \text{and} \quad \boldsymbol{\epsilon}_\theta^u(\mathbf{x}_t, t)
    \end{equation}

    \item Blend these predictions with a guidance scale parameter $\omega$:
    \begin{equation}
    \hat{\boldsymbol{\epsilon}}_\theta(\mathbf{x}_t, t, c) = \boldsymbol{\epsilon}_\theta^u(\mathbf{x}_t, t) + \omega \cdot (\boldsymbol{\epsilon}_\theta^c(\mathbf{x}_t, t, c) - \boldsymbol{\epsilon}_\theta^u(\mathbf{x}_t, t))
    \end{equation}

    \item Use this guided noise prediction in the reverse diffusion process:
    \begin{equation}
    \mathbf{x}_{t-1} = \frac{1}{\sqrt{\alpha_t}} \left( \mathbf{x}_t - \frac{1-\alpha_t}{\sqrt{1-\bar{\alpha}_t}} \hat{\boldsymbol{\epsilon}}_\theta(\mathbf{x}_t, t, c) \right) + \sigma_t \mathbf{z}
    \end{equation}
\end{enumerate}

The parameter $\omega$ controls the strength of the class conditioning, with higher values pushing the generation more strongly toward the target class characteristics.

\subsection{Complete Synthetic Data Generation Algorithm}

We integrate all components into a comprehensive algorithm for generating synthetic time series data:

\begin{algorithm}
\caption{Guided Synthetic Time Series Generation}\label{alg:architecture}
\begin{algorithmic}[1]
\Statex \textbf{Input:} Original imbalanced dataset $\mathcal{D} = \{(\mathbf{x}_i, y_i)\}_{i=1}^N$, trained classifier $f_\phi$, trained class-conditional diffusion model $\boldsymbol{\epsilon}_\theta$, set of minority classes $\mathcal{C}_{minority}$, number of synthetic samples per target region $M$
\Statex \textbf{Output:} Synthetic dataset $\mathcal{D}_{syn}$
\State Extract embeddings $\mathbf{z}_i = E_\phi(\mathbf{x}_i)$ for all $\mathbf{x}_i \in \mathcal{D}$
\State $\mathcal{D}_{syn} \leftarrow \emptyset$
\For{each $c \in \mathcal{C}_{minority}$}
    \State Identify target regions $\{\mathcal{R}_1, \mathcal{R}_2, \ldots, \mathcal{R}_R\}$ for class $c$
    \For{each region $\mathcal{R}_r$}
        \State Select $K$ conditioning instances $\{\mathbf{x}_1^c, \mathbf{x}_2^c, \ldots, \mathbf{x}_K^c\} \subset \mathcal{R}_r$
        \For{$m = 1$ to $M$}
            \State Compute diffusion timestep $t' \sim \text{Uniform}(t_{min}, t_{max})$
            \State Apply partial forward diffusion to get $\{\mathbf{x}_{1,t'}, \mathbf{x}_{2,t'}, \ldots, \mathbf{x}_{K,t'}\}$
            \State Compute conditioning point $\bar{\mathbf{x}}_{t'} = \frac{1}{K} \sum_{i=1}^K \mathbf{x}_{i,t'}$
            \State Run reverse diffusion from $\bar{\mathbf{x}}_{t'}$ with class $c$ to generate $\hat{\mathbf{x}}_m$
            \State $\mathcal{D}_{syn} \leftarrow \mathcal{D}_{syn} \cup \{(\hat{\mathbf{x}}_m, c)\}$
        \EndFor
    \EndFor
\EndFor
\State \Return $\mathcal{D}_{syn}$
\end{algorithmic}
\end{algorithm}

This algorithm systematically generates synthetic samples for each minority class by targeting specific boundary regions where additional data would be most beneficial for improving classification performance.

In our implementation, we set $t=300$ as number of partial diffusion steps, $K=3$ conditioning instances per region, and $\omega=3$ for the classifier-free guidance scale. These hyperparameters can be tuned based on the specific dataset characteristics and the desired balance between diversity and locality of the generated samples.

\subsection{Implementation Details}

Our implementation utilizes PyTorch for all neural network components. The class-conditional U-Net architecture for time series follows the design shown in Algorithm \ref{alg:architecture}, with specialized 1D convolutions and self-attention mechanisms adapted for temporal data.

For the MIT-BIH arrhythmia dataset, we preprocess the ECG signals by segmenting them into fixed-length windows of 128 timesteps and normalizing each segment to the range $[-1, 1]$. The diffusion model is trained with a batch size of 14, using the Adam optimizer with a learning rate of $3 \times 10^{-4}$ for 100 epochs.

The forward diffusion process uses a linear noise schedule with $\beta_1 = 10^{-4}$ and $\beta_T = 0.02$ over $T=1000$ timesteps. For synthetic data generation, we identify boundary regions using an entropy threshold of $\gamma = 0.7$ and density estimation with $k=5$ nearest neighbors.

The implementation leverages an Exponential Moving Average (EMA) of model parameters with a decay rate of 0.995 to improve generation stability, and classifier-free guidance is applied with a scale factor of $\omega=3$ during the reverse diffusion process.

\section{Experiments}
This section details the experimental setup used to evaluate our proposed DCIDM framework. We describe the datasets, baseline methods for comparison, evaluation metrics, and the specific experimental protocol followed.

\subsection{Datasets}
% Describe the MIT-BIH dataset in more detail (classes, imbalance ratio, preprocessing if not fully covered in 3.7).
% Describe any other datasets used for evaluation.
... (Content to be added here: Detailed description of datasets used, including statistics like number of samples, features, classes, and imbalance ratios.) ...

\subsection{Baseline Methods}
To assess the effectiveness of DCIDM, we compare it against several baseline approaches:
% List and briefly describe each baseline.
% Examples:
% \begin{itemize}
%     \item No Augmentation: Training a classifier on the original imbalanced dataset.
%     \item SMOTE: ...
%     \item ADASYN: ...
%     \item Time Series GAN (e.g., TimeGAN): ...
%     \item Standard Conditional DDPM (without targeted instance guidance): ...
% \end{itemize}
... (Content to be added here: Description of baseline methods selected for comparison.) ...

\subsection{Evaluation Metrics}
Given the imbalanced nature of the datasets, we employ a suite of evaluation metrics that provide a comprehensive view of classifier performance, especially concerning minority classes:
% List and define metrics.
% Examples:
% \begin{itemize}
%     \item Precision, Recall, F1-score (macro, micro, and per-class, especially for minority classes)
%     \item Area Under the Precision-Recall Curve (AUC-PR)
%     \item Area Under the ROC Curve (AUC-ROC)
%     \item Geometric Mean (G-Mean)
%     \item Matthews Correlation Coefficient (MCC)
% \end{itemize}
... (Content to be added here: List and brief explanation of evaluation metrics used.) ...

\subsection{Experimental Protocol}
% Describe the classifier used for evaluation (architecture, training details).
% How was the data split (train/validation/test)? Cross-validation details?
% How many synthetic samples were generated? How were they incorporated into the training set?
% Any specific hyperparameter settings for the evaluation classifier or baselines.
... (Content to be added here: Detailed procedure of how experiments were conducted, including classifier details, data splitting, and augmentation strategy.) ...

\section{Results}
This section presents the empirical results of our experiments, comparing DCIDM against the baseline methods using the evaluation metrics defined previously.

\subsection{Quantitative Results}
% Present main results in tables.
% Compare DCIDM with baselines across all metrics and datasets.
% Use bolding to highlight best-performing methods.
% Report mean and standard deviation if multiple runs were performed.
... (Content to be added here: Tables summarizing performance metrics, statistical significance tests if applicable.) ...

\subsection{Qualitative Results (Optional but Recommended)}
% Show plots of generated time series vs. real samples.
% t-SNE or UMAP visualizations of embedding space before and after augmentation, or showing where synthetic samples land.
... (Content to be added here: Visualizations of generated data quality, embedding space changes, etc.) ...

\subsection{Ablation Studies}
To understand the contribution of each component of DCIDM, we conduct several ablation studies:
% Describe each ablation study and present its results.
% Example studies:
% \begin{itemize}
%     \item Impact of targeted region identification (vs. random minority oversampling with diffusion).
%     \item Impact of partially diffused conditioning (vs. simpler conditioning or no instance guidance).
%     \item Sensitivity to key hyperparameters (e.g., $t'$, $K$, $\omega$).
% \end{itemize}
... (Content to be added here: Results of experiments designed to test the impact of individual components of the proposed method.) ...

\section{Discussion}
In this section, we interpret the experimental results, discuss their implications, and relate them to our initial contributions. We also acknowledge the limitations of our approach and consider its broader impact.

\subsection{Interpretation of Results}
% Discuss why DCIDM performs well (or not, in some cases).
% Connect performance to the mechanisms described in the Methodology.
... (Content to be added here: Analysis of what the results mean and why they occurred.) ...

\subsection{Connection to Contributions}
% Revisit the contributions listed in the Introduction and explain how the results support them.
% Elaborate on the theoretical analysis (Contribution 5) in light of the empirical findings.
... (Content to be added here: Linking experimental outcomes back to the paper's stated contributions.) ...

\subsection{Advantages and Disadvantages}
% Summarize the key strengths of DCIDM.
% Honestly discuss any limitations, computational costs, or scenarios where it might be less effective.
... (Content to be added here: A balanced view of the method's strengths and weaknesses.) ...

\subsection{Broader Impact and Ethical Considerations}
% Discuss potential positive applications.
% Address any ethical concerns or potential misuse (e.g., fairness, privacy if generating sensitive data).
... (Content to be added here: Reflection on the wider implications of the research.) ...

\section{Conclusion}
% Summarize the paper: the problem, the proposed method (DCIDM), and its key features.
% Reiterate the main findings and their significance.
... (Content to be added here: A concise summary of the paper's core message and achievements.) ...




\bibliography{References}
\bibliographystyle{plain}


\end{document}